\scanpage{80}
\mspage{14}
paire ne change rien; car si on remplace $D'_i$ par $\lambda D'_i$, donc $D''_i$ par $1/\lambda\,D''_i$, avec $\lambda\in\mathbb R^{*+}$ voisin de $1$, les \unsure{mêmes|dites|trois} relations \unsure{vaudront|voudront} \unsure{pour|fixer} les disques obtenus dans \inserted{le normalisé} $X''$ (c'est probablement tout à fait faux !), et on obtiendrait la $=$ surface recollée $X$ mais les $\Gamma_i$ remplacés par des cercles ``voisins'' $\Gamma'_i$.

Le cas-clef à comprendre sans doute est celui où $(g,\nu)=(0,4)$, i.e. le mécanisme de recollement de deux ``pantalons'' de Thurston, donc pour \deleted{\unsure{un ens}} \inserted{deux} ens$'$ $\unsure{\mathscr E|E|\Sigma}',\mathscr E''$ à deux éléments, fixés, les \unsure{foncteurs|fonctions|fonction} des
\[
\left\{\begin{aligned}
&\text{pantalons de types combinatoires } E'\amalg\{s\},\ E''\amalg\{t\},\\
&\text{recollés suivant } \{s,t\},\ \text{\deleted{donné} suivant}\\
&\text{les données }(l_i)_{i\in E'\amalg E''\amalg u}\in(\mathbb R^{*+})^{E'\amalg E''\amalg u},\\
&\text{et }\theta\in\mathbb U_{E'\wedge E''}\text{ qui fixe le recollement}
\end{aligned}\right.
\]
vers
% archive sheet number: 82
